\documentclass[11pt]{article}
\usepackage[margin=1in]{geometry}
\usepackage{amsmath}
\usepackage{booktabs}
\usepackage{enumitem}
\usepackage{hyperref}
\usepackage{xcolor}
\usepackage{framed}

\title{\textbf{Data Quality Dashboard User Guide}}
\author{PruTech Task C - Customer Query Analytics}
\date{\today}

\begin{document}

\maketitle

\tableofcontents
\newpage

\section{Dashboard Overview}

The Data Quality Dashboard is your command center for maintaining high-quality annotation data. It helps you identify problematic samples, understand annotator performance, and strategically improve your dataset through targeted quality improvements.

\subsection{What This Dashboard Solves}
\begin{itemize}
\item \textbf{Quality Control}: Find the most problematic samples requiring review or re-annotation
\item \textbf{Annotator Performance}: Understand which annotators need additional training
\item \textbf{Strategic Improvement}: Build better "gold-set" reference data for ongoing quality assurance
\item \textbf{Cost Optimization}: Focus your quality efforts where they'll have the biggest impact
\end{itemize}

\subsection{Dashboard Structure}
The dashboard contains three specialized analysis tabs:
\begin{enumerate}
\item \textbf{Top Disagreement Items} - Identify samples with highest annotation disagreement
\item \textbf{Per-Annotator Confusion} - Analyze individual annotator performance vs. consensus
\item \textbf{Suggested Gold-Set Refresh} - Get strategic recommendations for quality improvement
\end{enumerate}

\subsection{Who Should Use This Dashboard}
\begin{itemize}
\item \textbf{QA Teams}: Identify samples requiring review and re-annotation
\item \textbf{Training Managers}: Understand which annotators need focused training
\item \textbf{Data Scientists}: Build better training datasets with strategic sample selection
\item \textbf{Project Managers}: Track annotation quality trends and make resource allocation decisions
\end{itemize}

\section{Getting Started}

\subsection{Basic Navigation}
\begin{itemize}
\item \textbf{Theme Toggle}: Light/Dark buttons in top-right corner
\item \textbf{Tab Selection}: Click any tab name to switch between analyses
\item \textbf{Configuration Sections}: Each tab has a blue configuration card with controls
\item \textbf{Calculate Buttons}: Always click the "Calculate" or "Generate" button to run analysis
\item \textbf{Results Sections}: Scroll down to see charts, tables, and insights after calculation
\end{itemize}

\subsection{Understanding Result Types}
\begin{itemize}
\item \textbf{Disagreement Scores}: 0.0 = perfect agreement, 1.0 = maximum disagreement
\item \textbf{Agreement Rates}: Percentage of annotators who chose the same label
\item \textbf{Accuracy Scores}: How often an annotator matches the majority vote
\item \textbf{Confidence Levels}: Statistical confidence in majority vote decisions
\end{itemize}

\section{Tab 1: Top Disagreement Items}

\subsection{Purpose}
This tab identifies the most problematic samples in your dataset - documents where annotators disagree the most. These are prime candidates for review, discussion, and potential re-annotation.

\subsection{When to Use This Tab}
\begin{itemize}
\item Weekly quality reviews to identify problematic samples
\item Before finalizing datasets for machine learning training
\item When you notice declining annotation quality and need to find root causes
\item To create discussion materials for annotator training sessions
\end{itemize}

\subsection{Controls and Configuration}

\subsubsection{Label Type Dropdown}
\begin{itemize}
\item \textbf{Full Hierarchical Labels}: Complete labels (e.g., "Technical\_Performance\_Issue")
\item \textbf{L1 (Parent) Labels}: High-level categories (e.g., "Technical", "Billing")
\item \textbf{L2 (Child) Labels}: Specific sub-types (e.g., "Performance\_Issue", "Refund\_Request")
\item \textbf{Recommendation}: Start with "Full Hierarchical" for comprehensive disagreement analysis
\end{itemize}

\subsubsection{Annotator Selection}
\begin{itemize}
\item \textbf{What it controls}: Which annotators to include in disagreement calculations
\item \textbf{Default}: All annotators selected
\item \textbf{When to adjust}:
\begin{itemize}
\item Remove specific annotators to see if they're causing disagreements
\item Compare disagreement patterns between different annotator groups
\item Focus on a core group of experienced annotators
\end{itemize}
\end{itemize}

\subsubsection{Top N Samples Dropdown}
\begin{itemize}
\item \textbf{Options}: 10, 25, 50, or 100 most disagreeable samples
\item \textbf{Recommendation}: Start with 25 for manageable review workload
\item \textbf{Use 50-100}: When conducting comprehensive quality audits
\item \textbf{Use 10}: For quick daily quality checks
\end{itemize}

\subsubsection{Disagreement Score Sliders}
\begin{itemize}
\item \textbf{Min Disagreement Score (0.0-1.0)}: Exclude samples with too little disagreement
\item \textbf{Max Disagreement Score (0.0-1.0)}: Exclude samples with extreme disagreement
\item \textbf{Typical Settings}:
\begin{itemize}
\item \textbf{Min: 0.3, Max: 1.0}: Focus on moderately to highly disagreeable samples
\item \textbf{Min: 0.0, Max: 0.7}: Exclude completely chaotic samples, focus on learnable disagreements
\item \textbf{Min: 0.5, Max: 1.0}: Only the most problematic samples
\end{itemize}
\end{itemize}

\subsection{Understanding the Results}

\subsubsection{Disagreement Patterns Summary Cards}
\begin{itemize}
\item \textbf{Total Documents}: Overall dataset size for context
\item \textbf{Perfect Agreement}: Documents where all annotators agreed (quality baseline)
\item \textbf{High Disagreement}: Documents requiring immediate attention
\item \textbf{Average Disagreement}: Overall quality health indicator
\item \textbf{Most Disagreeable Document}: Worst-case example for review
\end{itemize}

\subsubsection{Disagreement Distribution Histogram}
\begin{itemize}
\item \textbf{X-axis}: Disagreement score (0 = perfect agreement, 1 = maximum disagreement)
\item \textbf{Y-axis}: Number of documents with that disagreement level
\item \textbf{Ideal Shape}: Tall peak at left (many documents with low disagreement)
\item \textbf{Warning Signs}: Flat distribution or peaks toward the right side
\end{itemize}

\subsubsection{Label Confusion Matrix}
\begin{itemize}
\item \textbf{What it shows}: Which labels are most often confused with each other
\item \textbf{Diagonal}: Correct classifications (darker = more common)
\item \textbf{Off-diagonal}: Confusion patterns (darker = more confusion)
\item \textbf{Business insight}: Identify label pairs requiring clearer guidelines
\end{itemize}

\subsubsection{Text Complexity vs. Disagreement Scatter Plot}
\begin{itemize}
\item \textbf{X-axis}: Text length or word count
\item \textbf{Y-axis}: Disagreement score
\item \textbf{What to look for}:
\begin{itemize}
\item \textbf{Positive correlation}: Longer texts are harder to annotate consistently
\item \textbf{No correlation}: Text length doesn't affect annotation difficulty
\item \textbf{Outliers}: Short texts with high disagreement (need guideline clarification)
\end{itemize}
\end{itemize}

\subsubsection{Top Disagreement Table}
\begin{itemize}
\item \textbf{Columns explained}:
\begin{itemize}
\item \textbf{Rank}: Disagreement severity ranking (1 = most disagreeable)
\item \textbf{Sample ID}: Document identifier for tracking
\item \textbf{Disagreement Score}: Quantified disagreement level
\item \textbf{Agreement Rate}: Percentage of annotators who agreed
\item \textbf{Unique Labels}: Number of different labels assigned
\item \textbf{Annotator Labels}: What each annotator chose (detailed breakdown)
\item \textbf{Text Preview}: Sample content for context
\end{itemize}
\item \textbf{Sorting}: Click column headers to sort by different criteria
\item \textbf{Filtering}: Use the filter boxes to search for specific patterns
\end{itemize}

\subsection{Business Actions Based on Results}

\subsubsection{High Disagreement Samples (Score > 0.7)}
\textbf{Immediate Actions}:
\begin{itemize}
\item Schedule team discussion sessions using these examples
\item Review and clarify annotation guidelines
\item Consider if these samples represent edge cases requiring special handling
\item Re-annotate with expert annotators for gold-set inclusion
\end{itemize}

\subsubsection{Moderate Disagreement Samples (Score 0.3-0.7)}
\textbf{Recommended Actions}:
\begin{itemize}
\item Use for annotator training exercises
\item Identify common confusion patterns from the confusion matrix
\item Develop specific guidance for these challenging scenarios
\item Consider creating sub-categories for frequently confused labels
\end{itemize}

\subsubsection{Text Complexity Patterns}
\textbf{If longer texts have more disagreement}:
\begin{itemize}
\item Provide guidelines for handling complex, multi-issue texts
\item Train annotators to identify primary vs. secondary issues
\item Consider breaking down complex texts into smaller annotation units
\end{itemize}

\section{Tab 2: Per-Annotator Confusion}

\subsection{Purpose}
This tab analyzes individual annotator performance by comparing each annotator's choices against the majority vote consensus. It helps identify which annotators need targeted training and what specific mistakes they're making.

\subsection{When to Use This Tab}
\begin{itemize}
\item Monthly annotator performance reviews
\item Before making training resource allocation decisions
\item When investigating quality issues affecting specific annotators
\item To identify high-performing annotators for mentorship roles
\end{itemize}

\subsection{Controls and Configuration}

\subsubsection{Analysis Mode Dropdown}
\begin{itemize}
\item \textbf{Complete Analysis}: Run all performance analyses (recommended)
\item \textbf{Individual Performance Only}: Focus just on accuracy metrics
\item \textbf{Pairwise Comparison Only}: Analyze annotator-to-annotator agreement patterns
\item \textbf{Recommendation}: Use "Complete Analysis" for comprehensive review
\end{itemize}

\subsubsection{Label Type Dropdown}
Same options as Tab 1 - choose which level of annotation to analyze for performance.

\subsubsection{Annotator Selection}
Choose which annotators to include in the performance analysis.

\subsection{Understanding the Results}

\subsubsection{Annotator Performance Ranking}
\begin{itemize}
\item \textbf{What it shows}: Accuracy scores for each annotator vs. majority vote
\item \textbf{Ranking order}: Highest to lowest accuracy
\item \textbf{Accuracy interpretation}:
\begin{itemize}
\item \textbf{90\%$\leq$}: Excellent performance
\item \textbf{80-89\%}: Good performance  
\item \textbf{70-79\%}: Needs improvement
\item \textbf{<70\%}: Requires immediate training intervention
\end{itemize}
\end{itemize}

\subsubsection{Systematic Bias Analysis}
\begin{itemize}
\item \textbf{What it shows}: Common mistake patterns across all annotators
\item \textbf{Format}: "True Label → Annotator's Choice (Frequency)"
\item \textbf{Business insight}: Identify systematic guideline problems affecting multiple annotators
\end{itemize}

\subsubsection{Pairwise Agreement Matrix}
\begin{itemize}
\item \textbf{What it shows}: Agreement rates between each pair of annotators
\item \textbf{Diagonal}: Always 100\% (annotator agrees with themselves)
\item \textbf{How to read}: 
\begin{itemize}
\item High values (>80\%) = good agreement
\item Low values (<60\%) = significant disagreement
\item Patterns can reveal annotator groups or outliers
\end{itemize}
\end{itemize}

\subsubsection{Individual Annotator Dropdown}
\begin{itemize}
\item \textbf{Purpose}: Select specific annotator to see their detailed confusion matrix
\item \textbf{Matrix interpretation}:
\begin{itemize}
\item \textbf{Y-axis}: True labels (majority vote)
\item \textbf{X-axis}: Annotator's labels
\item \textbf{Diagonal}: Correct classifications (should be darkest)
\item \textbf{Off-diagonal}: Specific mistake patterns
\end{itemize}
\end{itemize}

\subsection{Business Actions Based on Results}

\subsubsection{Low-Performing Annotators (<80\% accuracy)}
\textbf{Immediate Actions}:
\begin{itemize}
\item Schedule one-on-one training sessions
\item Review their confusion matrix to identify specific problem areas
\item Provide targeted examples for their most common mistakes
\item Consider temporary quality review for all their annotations
\end{itemize}

\subsubsection{Systematic Bias Patterns}
\textbf{When multiple annotators make the same mistakes}:
\begin{itemize}
\item Review annotation guidelines for those specific label pairs
\item Clarify distinction criteria between confused labels
\item Add more examples to training materials
\item Consider if label definitions need revision
\end{itemize}

\subsubsection{Annotator Groups in Pairwise Matrix}
\textbf{If you see distinct clusters}:
\begin{itemize}
\item \textbf{Two high-agreement groups}: May indicate interpretation differences - need alignment
\item \textbf{One outlier annotator}: Individual training needed
\item \textbf{Generally low agreement}: Systematic guideline problems
\end{itemize}

\section{Tab 3: Suggested Gold-Set Refresh}

\subsection{Purpose}
This tab provides strategic recommendations for building a high-quality "gold-set" of reference annotations. Gold-sets are used for training, quality monitoring, and validating annotation consistency over time.

\subsection{When to Use This Tab}
\begin{itemize}
\item Quarterly gold-set reviews and updates
\item When launching new annotation projects
\item After major guideline changes require validation
\item To build training datasets for machine learning models
\end{itemize}

\subsection{Controls and Configuration}

\subsubsection{Label Type Dropdown}
Same as other tabs - choose which annotation level to optimize.

\subsubsection{Gold-Set Strategy Dropdown}
\begin{itemize}
\item \textbf{High Agreement Only}: Select only samples with strong annotator consensus
\begin{itemize}
\item \textbf{Use when}: Building training data or establishing clear examples
\item \textbf{Benefit}: High confidence in "correct" labels
\item \textbf{Limitation}: May miss edge cases and difficult scenarios
\end{itemize}
\item \textbf{Disagreement Focus}: Select samples with useful levels of disagreement
\begin{itemize}
\item \textbf{Use when}: Identifying challenging cases for guideline development
\item \textbf{Benefit}: Captures difficult edge cases requiring attention
\item \textbf{Limitation}: Requires expert review to establish "correct" labels
\end{itemize}
\item \textbf{Mixed Strategy}: Balanced combination of high-agreement and disagreement samples
\begin{itemize}
\item \textbf{Use when}: Building comprehensive gold-sets for ongoing quality monitoring
\item \textbf{Benefit}: Captures both clear examples and challenging edge cases
\item \textbf{Recommendation}: Best choice for most scenarios
\end{itemize}
\end{itemize}

\subsubsection{Samples Per Label Slider}
\begin{itemize}
\item \textbf{Range}: 1-10 samples per unique label
\item \textbf{Recommendation}: 3-5 samples per label for balanced coverage
\item \textbf{Higher values (7-10)}: When you need extensive coverage for complex labeling schemes
\item \textbf{Lower values (1-2)}: For quick gold-set creation or budget constraints
\end{itemize}

\subsubsection{High Agreement Threshold Slider}
\begin{itemize}
\item \textbf{Range}: 60\%-95\% agreement required
\item \textbf{Default}: 80\%
\item \textbf{Higher thresholds (90\%+)}: Very strict quality, fewer candidates
\item \textbf{Lower thresholds (70\%-80\%)}: More inclusive, larger candidate pool
\end{itemize}

\subsubsection{Disagreement Range Slider}
\begin{itemize}
\item \textbf{Purpose}: Define "useful disagreement" levels for challenging examples
\item \textbf{Range}: 20\%-80\% agreement
\item \textbf{Typical Setting}: 30\%-60\% (moderate disagreement, still learnable)
\item \textbf{Avoid}: Very low agreement (<20\%) indicates chaotic, un-learnable samples
\end{itemize}

\subsubsection{Annotator Selection}
Choose which annotators to include in the consensus calculations.

\subsection{Understanding the Results}

\subsubsection{Gold-Set Quality Overview Dashboard}
\begin{itemize}
\item \textbf{Coverage Metrics}: How well the recommended samples cover all label types
\item \textbf{Quality Scores}: Confidence levels in the recommended samples
\item \textbf{Balance Indicators}: Distribution across different agreement levels
\item \textbf{Total Recommendations}: Number of samples recommended for gold-set inclusion
\end{itemize}

\subsubsection{Label Coverage Comparison Chart}
\begin{itemize}
\item \textbf{What it shows}: How many samples exist vs. how many are recommended for each label
\item \textbf{Goal}: Roughly equal representation across all important labels
\item \textbf{Red flags}: Labels with zero or very few recommendations (coverage gaps)
\end{itemize}

\subsubsection{Candidate Distribution Chart}
\begin{itemize}
\item \textbf{High Confidence Candidates}: Samples with strong agreement (training examples)
\item \textbf{Disagreement Candidates}: Samples with useful disagreement (edge cases)
\item \textbf{Ideal balance}: Depends on strategy, but generally 60-70\% high confidence, 30-40\% disagreement focus
\end{itemize}

\subsubsection{Agreement Distribution Scatter Plot}
\begin{itemize}
\item \textbf{X-axis}: Document ID or sample index
\item \textbf{Y-axis}: Agreement level
\item \textbf{Colored regions}: Show your high agreement and disagreement thresholds
\item \textbf{Selected points}: Recommended samples for gold-set inclusion
\end{itemize}

\subsubsection{Recommended Gold-Set Candidates Table}
\begin{itemize}
\item \textbf{Columns}:
\begin{itemize}
\item \textbf{Sample ID}: Document identifier
\item \textbf{Label}: Consensus or majority label
\item \textbf{Agreement Level}: Percentage of annotators who agreed
\item \textbf{Confidence}: Statistical confidence in the consensus
\item \textbf{Strategy Type}: High confidence vs. disagreement focus
\item \textbf{Text Preview}: Sample content for context
\end{itemize}
\item \textbf{Actions}: Export this list for gold-set creation workflows
\end{itemize}

\subsubsection{Coverage Gap Analysis Table}
\begin{itemize}
\item \textbf{What it shows}: Labels that are under-represented in recommendations
\item \textbf{Why gaps occur}:
\begin{itemize}
\item Very rare labels with insufficient samples
\item Labels with consistently poor agreement
\item Labels that don't meet quality thresholds
\end{itemize}
\item \textbf{Actions}: Manual review needed for these labels
\end{itemize}

\subsection{Business Actions Based on Results}

\subsubsection{High-Quality Recommendations}
\textbf{For samples with >90\% agreement}:
\begin{itemize}
\item Include in training materials immediately
\item Use as "gold standard" examples in annotator training
\item Include in automated quality monitoring systems
\item Consider for machine learning training data
\end{itemize}

\subsubsection{Moderate Disagreement Samples}
\textbf{For samples with 30-60\% agreement}:
\begin{itemize}
\item Schedule expert review sessions to establish consensus
\item Use in advanced annotator training for edge cases
\item Develop specific guidelines for these challenging scenarios
\item Consider as candidates for label refinement discussions
\end{itemize}

\subsubsection{Coverage Gaps}
\textbf{For under-represented labels}:
\begin{itemize}
\item Prioritize collecting more examples of these labels
\item Review if these labels are too rare to be useful
\item Consider combining very rare labels into broader categories
\item Manually curate examples if necessary for critical business cases
\end{itemize}

\section{Best Practices for Data Quality Management}

\subsection{Weekly Quality Workflow}
\begin{enumerate}
\item \textbf{Monday}: Run Tab 1 (Disagreement Analysis) to identify problem samples
\item \textbf{Wednesday}: Review Tab 2 (Annotator Performance) for training needs
\item \textbf{Friday}: Plan weekend training sessions based on identified issues
\end{enumerate}

\subsection{Monthly Strategic Review}
\begin{enumerate}
\item \textbf{Week 1}: Run comprehensive disagreement analysis
\item \textbf{Week 2}: Conduct annotator performance reviews
\item \textbf{Week 3}: Update gold-set using Tab 3 recommendations
\item \textbf{Week 4}: Implement training improvements and guideline updates
\end{enumerate}

\subsection{Quality Metrics to Track}
\begin{itemize}
\item \textbf{Average disagreement score}: Should decrease over time
\item \textbf{Perfect agreement percentage}: Should increase over time  
\item \textbf{Annotator accuracy scores}: Should improve with training
\item \textbf{Gold-set coverage}: Should be comprehensive across all important labels
\end{itemize}

\subsection{Red Flags Requiring Immediate Attention}
\begin{itemize}
\item \textbf{Disagreement scores increasing week-over-week}
\item \textbf{Any annotator with <70\% accuracy consistently}
\item \textbf{Systematic bias patterns affecting multiple annotators}
\item \textbf{Large coverage gaps in gold-set recommendations}
\end{itemize}

\section{Integration with Business Processes}

\subsection{QA Team Integration}
\begin{itemize}
\item Use disagreement analysis results for weekly review prioritization
\item Export problem samples for detailed manual review
\item Track improvement metrics after implementing recommended changes
\end{itemize}

\subsection{Training Team Integration}
\begin{itemize}
\item Use annotator performance data for personalized training plans
\item Incorporate high-disagreement samples into training materials
\item Monitor training effectiveness through improved accuracy scores
\end{itemize}

\subsection{Data Science Team Integration}
\begin{itemize}
\item Use gold-set recommendations for machine learning training data
\item Validate model performance using high-confidence samples
\item Identify data collection priorities based on coverage gap analysis
\end{itemize}

\section{Troubleshooting}

\subsection{Common Issues}
\begin{itemize}
\item \textbf{No results after clicking Calculate}: Ensure annotators are selected and thresholds are reasonable
\item \textbf{Very few recommendations}: Lower quality thresholds or expand disagreement ranges
\item \textbf{Performance analysis showing all low scores}: Check if majority vote calculation is working correctly
\end{itemize}

\subsection{Data Quality Issues}
\begin{itemize}
\item \textbf{Extremely high disagreement across all samples}: Review annotation guidelines for clarity
\item \textbf{No coverage gaps but poor agreement}: Focus on annotator training rather than sample selection
\item \textbf{Perfect agreement everywhere}: May indicate insufficient annotation diversity or complexity
\end{itemize}

\section{Conclusion}

The Data Quality Dashboard provides comprehensive tools for maintaining and improving annotation quality. Regular use of all three tabs creates a complete quality management workflow:

\begin{itemize}
\item \textbf{Identify Problems}: Tab 1 finds the most challenging samples
\item \textbf{Understand Causes}: Tab 2 reveals annotator-specific issues  
\item \textbf{Strategic Improvement}: Tab 3 guides targeted quality investments
\end{itemize}

Success comes from consistent usage, data-driven decision making, and continuous improvement based on the insights provided. Use this dashboard as your central hub for all annotation quality management activities.

\end{document}